% =====================================================================
%  CH 1-4 : Vision · Couche I et taxonomie · Couches et traçabilité · Carte
% =====================================================================
\section{Vision et objectifs}\label{sec:vision}

Le développement logiciel est entré dans l’ère \emph{agent-first} : une part croissante du
code est écrite, relue et modifiée par des modèles de langage plutôt que par des humains.
Cette bascule remet en question une hypothèse centrale de nos outils — que le code source
est optimisé pour la lecture humaine. TRIADE est la réponse disciplinée à ce réexamen :
plutôt que de supprimer les langages (une impossibilité de principe, puisqu’un modèle de
langage manipule des symboles textuels et que tout format binaire rendu lisible par un
agent constitue de facto un langage), TRIADE déplace le langage vers une fine couche de
spécification maintenue par l’agent, tandis que l’artefact livré est du code machine natif
pur. Le compilateur disparaît du pipeline de livraison : ce que l’agent écrit est ce qui
s’exécute — et rien d’autre.

Le nom encode l’architecture : trois couches indissociables — Intention (I), Binaire (B),
Carte (C) — unies par un verrou croisé cryptographique. La couche I contient le quoi et le
pourquoi : contrats typés, pré et post-conditions, invariants, cas d’échec, vecteurs de
vérification. La couche B contient le comment : des octets scellés, versionnés, à adresses
stables. La couche C est l’index naviguable qui lie les deux et sert d’unique interface de
développement à l’agent.

La v0.4.0 tire les enseignements de la première validation à l’échelle — le recodage
complet de RoomEq (démonstration v0.3) — et les convertit en normes : ce que la
démonstration a pratiqué (reverse-specs, pré-versions, partage de fixtures, mesure avant
optimisation) devient exigence ; ce que son analyse des défauts a révélé devient
taxonomie ; ce que son étude de cas établit et ne prouve pas reste explicitement hors
périmètre de preuve.

La v0.5.0 est la révision du retour du terrain. Les jalons M0 (chaîne artefact), M1
(mini-libc scellée) et M2 (patch à chaud en production, avec retour arrière) sont
franchis — chacun sur validation terrain Windows 11, chacun scellé par une balise
finale immuable (0.2.0, 0.3.0). La campagne M2 a coûté huit classes de défauts
invisibles à l’émulation (§\ref{sec:taxonomieX}) — le prix d’une leçon : un oracle
valide un artefact contre un modèle d’environnement, et la complétude du modèle ne se
démontre qu’en exécutant le binaire sur le système réel. Ce que le terrain a appris
devient norme : la boucle d’évidence (§\ref{sec:terrain}), la taxonomie X,
l’instrumentation contractualisée — et l’ouverture du jalon M3, contrats renforcés
sur entrées non fiables.

\subsection{Fonction objectif hiérarchisée — invariant de révision}\label{sec:objectif}

\begin{center}
\large\bfseries Sécurité $>$ Performance $>$ Optimisation $>$ Coût en tokens
\end{center}

\medskip

Arbitrage lexicographique : un objectif n’est jamais sacrifié au profit d’un objectif de
rang inférieur. La v0.4.0 consigne cette hiérarchie comme \textbf{invariant de révision}
(D19) : aucune révision future du protocole ne peut la modifier sans invalider
explicitement le protocole entier — pas l’ordre, pas les rangs, pas la non-compensation.

\begin{table}[htbp]
\centering\small
\caption{La fonction objectif et ses métriques de validation.}
\label{tab:objectif}
\begin{tabularx}{\textwidth}{@{}l Y Y@{}}
\toprule
\textbf{Objectif} & \textbf{Signification concrète} & \textbf{Métrique de validation} \\
\midrule
Sécurité & Sûreté mémoire validée par oracle ; intégrité carte $\leftrightarrow$ binaire ;
durcissement au chargement & 0 corruption détectée ; tout échec de verrou bloque la
session ; échec au lancement = arrêt propre \\
Performance & 100 \% code machine natif ; zéro abstraction intermédiaire &
Comparaison par bloc scellé face à une référence compilée \\
Optimisation & Descente à l’ISA sur hot paths identifiés — jamais sans mesure (R11) &
Gain mesuré et journalisé par bloc optimisé \\
Coût en tokens & Édition ciblée via la carte ; jamais de lecture intégrale du binaire &
Budgets par type d’opération (§\ref{sec:carte}) \\
\bottomrule
\end{tabularx}
\end{table}

\subsection{Périmètre et non-objectifs}\label{sec:perimetre}

Cinq non-objectifs. Premièrement, TRIADE ne prétend pas remplacer les langages dans
l’industrie : c’est une architecture de laboratoire. Deuxièmement, la furtivité face aux
EDR n’est pas recherchée : l’outil cible est un utilitaire de pentest pour tests autorisés
en laboratoire, appelant les API du système via les interfaces standards du profil.
Troisièmement, l’universalité est structurelle mais progressive : chaque profil doit être
instancié et validé (M6, M7) avant d’être déclaré supporté — le protocole ne promet pas la
portabilité, il la vérifie. Quatrièmement, aucune couche d’abstraction à l’exécution
(bytecode portable, machine virtuelle, traduction à la volée) n’est introduite : elle
trahirait la fonction objectif Performance. Cinquièmement, la couche de gouvernance
validée par la démonstration et le cœur natif du protocole (notation, verrou
cryptographique sur octets, patch à chaud) restent des niveaux de preuve distincts : les
jalons M0 à M7 demeurent la charge de validation du cœur natif.

\subsection{Ce que change la v0.4.0}\label{sec:nouveautes}

Dix-sept évolutions, toutes tracées dans la correspondance du §\ref{sec:correspondance} et
dans le journal des versions, toutes rattachées aux décisions D19 à D24 : normalisation de
la couche I, taxonomie des défauts, traçabilité I $\leftrightarrow$ C $\leftrightarrow$
B $\leftrightarrow$ tests, norme de fragmentation, localisation par intention, R10
renforcée, R11 nouvelle, suite de conformité par profil, parité inter-profils,
pré-versions, reverse-specs, journal d’optimisation, modèle de menace étendu, journal des
décisions normalisé, tableau de bord de conformité. Aucune règle existante n’est
affaiblie ; aucune couche n’est ajoutée à l’exécution ; la hiérarchie de la fonction
objectif est inchangée.

\subsection{Ce que change la v0.5.0}

Huit évolutions, toutes tracées dans la correspondance (§\ref{sec:evolutions050}) et le
journal des versions, toutes rattachées aux décisions D26 à D45 : la validation
terrain devient une couche de vérification normée — boucle d’évidence, journal
d’exécution, instrumentation contractualisée, diff chirurgical (§\ref{sec:terrain}) ;
la taxonomie des défauts s’étend des classes T (contrat non écrit) aux classes X —
défauts invisibles à l’émulation, huit classes nommées par la campagne M2
(§\ref{sec:taxonomieX}) ; l’état des jalons est constaté — M0, M1 et M2 franchis sur
validations terrain Windows 11, balises finales 0.2.0 et 0.3.0 posées
(§\ref{sec:jalons}) ; le budget de localisation est corrigé à 64 tokens (D31) et les
budgets affichent leurs valeurs mesurées ; les volumes d’oracle effectifs sont
consignés ; les preuves terrain sont archivées brutes avec leurs empreintes ; le
corpus de démonstration (huit itérations instrumentées) devient l’occurrence de
référence de la méthode ; le jalon M3 est ouvert — contrats renforcés sur entrées non
fiables et fuzzing à propriétés (D45). Aucune règle n’est modifiée ; aucune couche
n’est ajoutée à l’exécution ; la hiérarchie de la fonction objectif est inchangée.

\subsection{Ce que change la v0.5.1}\label{sec:nouveautes051}

Un composant nouveau et deux constats, tous tracés dans la correspondance
(§\ref{sec:evolutions051}) et le journal des versions, rattachés à la décision D75 : le
\textbf{coffre-fort des phrases maîtresses} — conteneur chiffré TRIADE-CF1 centralisant
toutes les phrases maîtresses par projet, sur la pile cryptographique de la carte au
repos (Argon2id aux planchers, HKDF à contexte dédié, XChaCha20-Poly1305), doté d’un
autotest de falsification et de règles d’export/import (continuité,
§\ref{sec:gouvernance} ; FORMAT\_CARTE § 11) ; l’état des jalons est constaté — M3
franchi le 28 octobre 2026 sur la campagne million (D72), la conformité permanente
(D73) et le scellage de la balise finale 0.4.0 (D74) ; le modèle de menace gagne la
ligne du vol du coffre-fort (§\ref{sec:menaces}). Aucune règle n’est modifiée ; aucune
couche n’est ajoutée à l’exécution ; la hiérarchie de la fonction objectif est inchangée.

\subsection{Ce que change la v0.5.2}\label{sec:nouveautes052}

Une évolution de discipline, tracée dans la correspondance (§\ref{sec:evolutions052}) et
le journal des versions, rattachée à la décision D81 : la \textbf{phrase de contrôle du
coffre-fort passe sous garde durable} — dès sa production (init, rotation,
changement), elle est sauvegardée dans le dépôt git (commit dédié, push GitHub, miroir
du coffre Drive par le hook) et au racine du coffre Drive \emph{avant} toute autre action
de la session ; le fichier \code{PHRASE-CONTROLE-COFFRE.txt} (phrase seule sur une
ligne, lisible tel quel par \code{--phrase-fichier}) \emph{est} la remise — l’affichage
unique de cérémonie (D27/D30) cesse d’être le canal de survie. Le constat fondateur : la
rotation D80 avait produit une phrase neuve dont la remise prévue en fin de tour n’a
jamais eu lieu — la session d’exécution est morte avant (pannes d’infrastructure
récurrentes) — et la phrase ne vivait plus que dans le bac à sable éphémère : un décès
d’infra de plus et le coffre entier devenait indéchiffrable à jamais. L’arbitrage est
consigné : le risque de \textsc{perte} (vécu, répété) est jugé supérieur au risque
d’\textsc{exposition} (URL non devinable) ; le modèle de menace en tire la conséquence
honnête — la confidentialité des phrases au repos repose désormais sur le contrôle
d’accès au dossier du coffre Drive, la mitigation première devient la restriction du
partage par lien, et une phrase entrée dans l’historique git y demeure à vie : on ne
compte plus sur sa confidentialité, la rotation reste le remède. L’état des jalons est
constaté — M4 ouvert le 28 octobre 2026 (D77, couche I écrite-avant-code : huit
spécifications \code{tls} et deux révisions \code{http}). Aucune règle n’est modifiée ;
aucune couche n’est ajoutée à l’exécution ; la hiérarchie de la fonction objectif est
inchangée.

\subsection{Ce que change la v0.6.0}\label{sec:nouveautes060}

La révision de génération attendue par le plan : le \textbf{programme de souveraineté
cryptographique est accompli et prouvé sur terrain} — décision de passation D133,
arbitrage du développeur. Trois axes livrés (§\ref{sec:souverainete}) : l'\textbf{axe 2}
(D86) — le coffre CF2 à clé maîtresse répartie par Shamir, cinq gardes seuil trois,
sentinelle à double seuil puis PORTE D120 ; l'\textbf{axe 3} (D87) — le sceau de
l'Intention : dix arbres de couche I signés par la clé Ed25519 hors-ligne du
développeur, gardes de build et de patch ; l'\textbf{axe 4} (D121, né du terrain) — le
coffre CF3 : la serrure à triple scellement, le piège sur média séparé, la sentinelle
instantanée \emph{par invocation} et le cycle en deux actes. L'\textbf{axe 1} (D85,
attestation TEE) demeure une exigence consignée : l'infrastructure sans racine de
confiance matérielle ne peut pas la satisfaire et le dit. La \textbf{chaîne réelle du
coffre} a migré sur le terrain du développeur les 30--31 octobre 2026 — CF1 vérifié
INTÈGRE, phrase de contrôle changée (D110), CF2 migré (cinq gardes, seuil 3, trois
copies hors-ligne), CF3 migré (époque 1, checksum \code{0cc89a96}), verdict à froid
VERT deux fois — les cérémonies étant conduites par le maître de cérémonie graphique
(D107--D118) avec le kit Drive tenu identique au miroir par construction (D130). Un
\textbf{coffre neuf naissait CF3} : \code{init} direct à l'époque 1 — la chaîne
de migration est l'histoire du coffre réel, pas un rite obligé (la v0.6.1 remet ce
choix au développeur — §\ref{sec:nouveautes061}, §\ref{sec:initniveau}). Un document
complémentaire, le \emph{Livre du Coffre} (même date, même dépôt), porte la posture
opératoire du maître : cérémonies pas à pas, runbook, mode dégradé. Aucune règle du
cœur n'est modifiée ; aucune couche n'est ajoutée à l'exécution ; la hiérarchie de la
fonction objectif est inchangée.

\subsection{Ce que change la v0.6.1}\label{sec:nouveautes061}

Une seule règle, explicitement invalidée et remplacée (§\ref{sec:evolutions061}) :
le niveau de sécurité d'un coffre neuf. La v0.6.0 (D133) l'avait fixé — un coffre
neuf naît CF3, \code{init} direct à l'époque 1 ; la v0.6.1 (D137) en fait
\textbf{l'arbitrage du développeur au moment de créer son projet de développement} :
CF1, CF2 ou CF3, chacun avec son acte de naissance, ses obligations et son guide
(§\ref{sec:initniveau}) — imposer CF3 par défaut risque d'être trop contraignant
selon le projet : média du piège à produire dès le premier jour, cycle en deux actes
à chaque session, sachet physique à tenir — la contrainte se justifie quand le secret
la mérite. La chaîne de migration CF1 $\to$ CF2 $\to$ CF3 demeure disponible à tout
moment pour un coffre dont les besoins grandissent ; le coffre réel de TRIADE reste
CF3 (époque 1, checksum \code{0cc89a96}). Aucune autre règle n'est modifiée ; aucune
couche n'est ajoutée à l'exécution ; la hiérarchie de la fonction objectif est
inchangée.

\subsection{Ce que change la v0.7.0}\label{sec:nouveautes070}

La révision de consolidation du programme d'après M7 (D160) : ce que la clôture de la
v0.6.1 tenait « hors périmètre courant » entre dans la spécification, sur résultats
mesurés — la clause est levée NOMMÉMENT pour les deux premières familles, les profils
additionnels demeurent hors périmètre. Le \textbf{jalon M8} (le bloc optionnel de
syscalls directs — le domaine \code{sys}) est franchi : la charte à huit règles (le SSN
est une donnée, jamais un octet de code ; le motif du stub vérifié avant toute
extraction ; un stub divergent est un refus nommé — le résolveur ne devine jamais, il ne
s'adapte jamais en silence), les spécifications sys-01…03 écrites-avant-code, le harnais
différentiel 61 contrôles verts sur fixtures \code{ntdll} synthétiques (le laboratoire
ne connaît jamais le SSN du terrain), la cérémonie de l'ère passée avant le build (12
arbres, D162/D163), le run terrain sur le VRAI \code{ntdll} — la twin de la ligne
hello-02 par les DEUX chemins (l'IAT \code{kernel32} et le syscall direct
\code{NtWriteFile}), le refus $-4$ nommé du stub crocheté — recoupé 30/30 sur pièces
croisées, la finale \code{0.9.0} scellée (Entrée 164, D165 : le scellage exige le
verdict du recoup comme pièce ; le scellage est un fait de carte, l'artefact ne bouge
jamais). Le \textbf{jalon M9} (TLS 1.3 pur en blocs scellés — le domaine \code{tls13},
l'ère 1.0.0) est ouvert (D166) : la couche I complète — 28 spécifications en six
familles (sha, hkdf, x25519, aead, enregistrement, poignée, principal), l'instrument
\code{verifier\_specs\_m9a.py} 117 contrôles verts (D167) — et la cérémonie de l'ère
passée le 05/11 : la remise du développeur, TREIZE arbres signés hors ligne à la clé du
04/11 (la continuité prouvée par le déterminisme Ed25519, D163 : les onze arbres
inchangés rendent les signatures de l'ère M8), le garde repassé VERT 13/13 — la dette
D166 soldée sur pièces. Le \textbf{sceau de l'Intention} porte désormais treize arbres
(§\ref{sec:sceauI}) et les cérémonies d'ère entrent au protocole
(§\ref{sec:ceremoniesere}) : la dette se nomme à la naissance, jamais après coup (D166,
règle 8) ; la charge se prépare au geste qui sert la main du développeur (D168). Aucune
règle du cœur n'est modifiée ; aucune couche n'est ajoutée à l'exécution ; la hiérarchie
de la fonction objectif est inchangée.

\subsection{Ce que change la v0.8.0}\label{sec:nouveautes080}

La révision de consolidation du programme d'après M9 (D170/D171). Le \textbf{jalon M9}
(TLS 1.3 pur en blocs scellés — le domaine \code{tls13}, l'ère 1.0.0) est franchi : les
28 blocs notés-mesurés-prouvés — les primitives aux vecteurs officiels octet pour octet
(FIPS 180-4, RFC 2104/4231, 5869, 7748, SP 800-38D), la poignée rejouée sur le
transcript RFC 8448, les héritages PAR LA CARTE (R7) —, le socle socket par syscalls
directs, le septième run terrain vert sur pièces et recoupé 34/34, LA SESSION RÉELLE
(la ligne applicative chiffrée ÉCHANGÉE avec un serveur réel, le témoin
« HTTP/1.1 200 OK » DÉCHIFFRÉ PAR LA MACHINE, l'épinglage $-6$ NOMMÉ contre le vrai
serveur, la fermeture propre), la finale \code{1.0.0} scellée sur le recoup exigé
comme pièce (D170 : le scellage est un fait de carte, l'artefact ne bouge jamais) — et
l'ère M8 refermée propre par l'export de continuité. \textbf{LE TEST GRANDEUR NATURE
est ouvert (D171)} : le protocole avait prouvé chaque capacité, M0 à M9 — il reste à
prouver qu'il DÉVELOPPE une application réelle NEUVE, de l'intention à l'octet livré
(le modèle RoomEq s'arrêtait au recodage ; le TGN exige la CRÉATION). L'application —
arbitrée à D171, contredisable — est « Loupe », l'éditeur hexadécimal au panneau
d'empreintes sur le profil primaire (l'outil du quotidien du projet : voir les octets,
leurs lignes, leurs empreintes md5/sha-256, leur place) : le domaine \code{loupe}
fondé (la charte + le tableau des 18 blocs en cinq familles), la couche I posée
écrite-avant-code (D172 : les 18 spécifications, le plan d'adressage arrêté avant
toute notation, la table des sept verdicts), et les QUATRE chantiers de notation
accomplis — le noyau fichier + empreintes (D173 : les héritages vivent à leurs
adresses scellées, le mini-OS virtuel du harnais), le rendu + la navigation (D174 :
l'oracle à scénarios — l'écran attendu octet par octet —, le mini-GUI virtuel, le
plancher qui ne recule jamais), l'édition + la table de routage étendue de loupe-09
r2 (D175 : le disque saisissable et la mort du harnais, le réseau de la garde
16$\leftrightarrow$17$\leftrightarrow$18, l'écriture à l'octet prouvée différentielle) —
les 18 blocs NOTÉS-MESURÉS-PROUVÉS, 45 + 53 + 59 contrôles verts aux trois graines sur
les trois harnais. \textbf{L'ère du TGN est OUVERTE AU BUILD} (v0.8.1, la révision de
terrain) : LA CÉRÉMONIE D'ÈRE est PASSÉE SUR PIÈCES — le garde à QUATORZE arbres
verts, la dette D166/D171 SOLDÉE, la remise du développeur INSTALLÉE au coffre
(§\ref{sec:ceremoniesere}) — le build \code{loupe.exe 0.1.0-dev.1} est NÉ, la carte
génération 13, le verrou croisé 26/26. \textbf{ET LE TERRAIN A DIT TROIS FOIS,
CHAQUE REFUS NOMMÉ} (v0.8.2) : le run M10-(e) a REFUSÉ dev.1 par le VRAI
chargeur Windows — 0xc000007b au geste de cartographie (LoadLibraryEx
DONT\_RESOLVE : l'image ASLR sans table de relocalisation, l'oubli nommé du
premier build) — puis dev.2 DU MÊME REFUS (le BASERELOC rendu était NÉCESSAIRE
mais PAS SUFFISANT) : la preuve par la lignée isole les axes par lesquels
l'image sortait de l'enveloppe des HUIT artefacts passés au terrain (le trou de page
0x3000-0x4000, le BSS massif \code{.data} VS~$\gg$~RS, le sous-système GUI) ;
LA RETOUCHE dev.3 ramène l'image DANS l'enveloppe prouvée (.rdata étendue 0x2000 :
plus aucun trou ; .data VS~$=$~RS : plus aucun BSS — AUCUNE adresse scellée ne
bouge), ET LE KIT PORTE SA CASCADE DIAGNOSTIQUE. \textbf{LE RUN REJOUÉ A ÉTÉ LE
PAS DÉCISIF (l'Entrée 206) : LE VRAI CHARGEUR DE WINDOWS A CARTOGRAPHIÉ L'IMAGE —
la première fois au terrain, l'enveloppe prouvée DEMEURE NORME DE FAIT — puis LE
TÉMOIN A NOMMÉ AU T0 [CADRE-NÉ]} : le CADRE \code{.data} n'était pas né poignée
INVALIDE, l'initialiseur du build avait posé le sentinel à +8 (UNE TAILLE
\code{CADRE\_POIGNEE} prise pour un offset) au lieu de +528 (la chaîne des
tailles du plan — le zéro aurait été UNE SESSION FANTÔME au premier \code{ouvrir()}) ;
LA RETOUCHE dev.4 fait de l'initialiseur \code{.data} LE MIROIR EXACT du
\code{CADRE\_VIERGE} du harnais (l'offset CALCULÉ, jamais gravé de tête) ET LE
GARDE EST GUÉRI : le validateur du build validait à +8 — la même constante
fautive, un garde blessé par le défaut qu'il croyait garder ne garde rien — il
relit désormais LE MÊME miroir (la comparaison ENTIÈRE, le \code{.data} relu à
l'octet près) et le scénario relit LE \code{.data} ENTIER contre le miroir — un
échec nommé vaut plus qu'un vert muet, et le TGN fait exactement ce pour quoi il
existe (§\ref{sec:terrain}). \textbf{LE QUATRIÈME RUN A ÉTÉ LE PAS LE
PLUS LOIN JAMAIS VÊCU — ET LA VITRINE A VÉCU SA PREMIÈRE FOIS} (v0.8.3,
l'Entrée 207) : \textbf{[T0] TOUT VERT au témoin} (l'image cartographiée, LE
CADRE NÉ MIROIR — la retouche de l'Entrée 206 TENUE au terrain, l'IAT 37
feuilles posées) PUIS \textbf{[T1] ÉCHEC NOMMÉ [SESSION]} : loupe-01 a rendu
$-6$ SYSTÈME sur un fichier pourtant relu à l'instant, ET « quand je lance
l'exe, la fenêtre s'ouvre et se referme immédiatement » — LA MORT SANS TÉMOIN.
L'AUTOPSIE SUR LA PIÈCE NOMME \textbf{LA CLASSE DES ARGUMENTS PILE} : les
arguments 5+ de l'ABI msx64 vivaient AU HOME SPACE \code{[rsp..0x18]} au lieu
de \code{[rsp+0x20]+} — le VRAI CreateFileW lisait sa disposition À LA PILE NON
INITIALISÉE (loupe-01) ET le contexte condensé VIVAIT SUR LE VRAI SLOT (loupe-05 :
le VRAI ReadFile lisait H0[0] \code{0x6A09E667…} COMME UN POINTEUR — la mort de
la vitrine avec fichier ; loupe-06 pareil) — POURQUOI LE LABO N'AVAIT RIEN VU :
les shims du harnais ne lisaient AUCUN argument pile (le garde avait déjà la
bonne convention dans le shim TextOutW — elle n'était pas appliquée aux feuilles
du noyau). LA RETOUCHE dev.5 : les TROIS notations à leurs vrais slots, DEUX AXES
D'ALIGNEMENT nommés par L'AUDIT D'ABI msx64 NÉ AU BUILD (loupe-10 le
sous-programme \code{.peindre} appelait TextOutW pile $\equiv 8$ : le VRAI
gdi32 est SSE — l'axe que seul l'audit statique voit ; loupe-11 le cadre
\code{0x38}$\to$\code{0x40}), LE GARDE GUÉRI PAR LA CONVENTION (les shims
relisent LES MÊMES SLOTS QUE LE VRAI — la falsification demandée est attrapée)
ET LE KIT PORTE DÉSORMAIS SON TÉMOIN VITRINE (\code{temoin-vitrine.bat} : les
deux codes de sortie + le rapport WER — la prochaine mort sera nommée).
\textbf{LE CINQUIÈME RUN A NOMMÉ $-7$ VERROUILLÉ — ET LA PIÈCE EST
ACQUITTÉE} (v0.8.4, l'Entrée 208) : l'appel CreateFileW de loupe-01 était
NÉ PARFAIT dans la pièce (rcx = le chemin, rdx = \code{0xC0000000}
[GENERIC\_READ|GENERIC\_WRITE], r8 = 1 [FILE\_SHARE\_READ], les arguments
5/6/7 À LEURS VRAIS SLOTS, la pile $\equiv 0 \pmod{16}$) — LA PREUVE PAR LE
NOM CHANGÉ : dev.4 rendait $-6$ SYSTÈME (gle 87 : l'appel mal né),
dev.5 rend $-7$ VERROUILLÉ (gle 32 : l'appel NÉ, refusé pour une RAISON
RÉELLE) — la classe de l'Entrée 207 A TENU AU TERRAIN ; ET LA RAISON
RÉELLE ÉTAIT LE KIT LUI-MÊME : le scénario EST LE CHARGEUR —
\code{LoadLibraryEx} DONT\_RESOLVE TIENT loupe.exe comme SECTION VIVANTE
(sans partage d'écriture), et le VRAI Windows refuse gle 32 d'ouvrir À
L'ÉCRITURE le fichier d'une image mappée : T1 offrait LA PIÈCE ELLE-MÊME
à l'écriture, IMPOSSIBLE PAR CONSTRUCTION — invisible aux runs
précédents parce qu'aucun appel n'était jamais assez bien né pour
atteindre le partage ; loupe-01 A TRADUIT LE REFUS AU NOM PRÈS (I4/E2 :
gle 32 $\to$ $-7$) — LA VOIE NOMMÉE DE L'ÈRE A VÉCU SON PREMIER TERRAIN,
VERTE PAR LE REFUS. LA RETOUCHE dev.6 EST AU KIT SEUL — L'ARTEFACT
INCHANGÉ (\code{loupe.exe 0.1.0-dev.5}, sha256 \code{6f830f9c…}) :
L'OMBRE (la session T1 naît sur \code{temoin-session.bin}, la copie
byte-à-byte de la pièce, née et morte au run — même loi que la sonde 205
et le témoin vitrine), LE FRÈRE DU DÉFAUT CHASSÉ (la session double de T6
offrait elle aussi la pièce sous le chargeur : l'offrande est l'ombre,
le premier appel vérifié), ET LA VITRINE SE VIT PAR SON TÉMOIN
(\code{temoin-vitrine.txt} REQUIS au réceptacle : une mort sans témoin
ne recoupe pas). Le projet
fils COFFRE-FORT est dégelé en parallèle
(D135 amendée par D171) : le père continue sans ralentir, la frontière fondatrice
demeure INTACTE. Aucune règle du cœur n'est modifiée ; aucune couche n'est ajoutée
à l'exécution ; la hiérarchie de la fonction objectif est inchangée.

\section{La couche I : format normalisé des spécifications}\label{sec:coucheI}

La v0.3 décrivait la couche I comme « contrats légers » (D3), sans en fixer le gabarit —
la légèreté demeure, la structure devient normative (D19). Une spécification I n’est pas un
document libre : c’est un artefact diffable dont la structure conditionne la vérification.

\subsection{Identifiants stables}

Chaque spécification reçoit à sa création un identifiant \code{<DOM>-<NN>} : un domaine
parmi une liste figée par projet et un numéro séquentiel. L’identifiant est immuable et
jamais réutilisé ; un module supprimé laisse son identifiant vacant à jamais. C’est
l’équivalent exact, côté intention, de l’adresse stable côté binaire (R7) : les références
croisées — dépendances, carte, tests — ne cassent jamais. Le couple (identifiant,
révision) est la clé de traçabilité du §\ref{sec:traceabilite}.

\subsection{Gabarit obligatoire}

\begin{lstlisting}
<DOM>-<NN> . <nom court> . revision <R>
Intention : <une phrase — le pourquoi>
Contrat :
  pre  : P1..Pn    (preconditions numerotees, une phrase chacune)
  post : Q1..Qn    (postconditions numerotees)
  inv  : I1..In    (invariants)
Echecs :
  E1 : <condition> -> <sanction nommee : exception | arret propre | repli>
Vecteurs :
  V1 : <vecteur dore | propriete differentielle | test de rejet>
Fixtures    : <references aux fixtures contractuelles partagees>
Provenance  : ecrite-avant-code | reverse-spec de <binaire scelle @ tag>
Decisions   : <Dxx...>
\end{lstlisting}

Sept blocs, tous obligatoires — les échecs et vecteurs ne peuvent être vides que pour un
module sans effet observable, situation rare à justifier dans le champ Intention. Trois
règles de rédaction conditionnent la diffabilité exigée par R2 : \textbf{une exigence par
phrase numérotée} (le diff montre exactement quelle exigence change) ; \textbf{aucune
connaissance dupliquée} (renvoi par identifiant, sinon défaut T5) ; \textbf{chaque
pre/post/invariant relie au moins un vecteur} (une exigence invérifiable est une exigence
mal formulée). La provenance distingue les spécifications écrites avant le code — le cas
normal, exigé par R2 — des reverse-specs (§\ref{sec:reversespecs}).

\subsection{Révision et cycle de vie}

La révision s’incrémente à chaque changement du contrat. Une évolution qui rompt un
contrat existant ne modifie pas la spécification — elle en crée une nouvelle, l’ancienne
étant dépréciée avec renvoi : miroir exact du cycle de vie des blocs. Les appelants
dépendent du couple (identifiant, révision) ; la carte porte ce couple dans chaque entrée.

\subsection{Fixtures contractuelles partagées}

Quand plusieurs consommateurs se lient au même contrat externe — sortie d’un outil
système, format de fichier, API distante — les échantillons qui définissent ce contrat
appartiennent à la spécification, pas aux consommateurs. Le champ Fixtures référence un
jeu unique ; toute copie privée est un défaut de classe T5. Le mécanisme est né dans la
démonstration RoomEq : l’installateur scellé et l’application recodée partagent les
fixtures de parsing de pnputil — le défaut historique qui a survécu à trois versions
vivait précisément de la duplication.

\subsection{Taxonomie des défauts de contrat non écrit}\label{sec:taxonomie}

La classe de défauts que le protocole existe pour éliminer a un nom : \emph{contrat non
écrit}. La v0.4.0 en fixe la taxonomie (D19), née de l’analyse des défauts historiques de
la démonstration.

\begin{table}[htbp]
\centering\small
\caption{Taxonomie des défauts de contrat non écrit.}
\label{tab:taxonomie}
\begin{tabularx}{\textwidth}{@{}l Y Y@{}}
\toprule
\textbf{Classe} & \textbf{Définition} & \textbf{Verrou} \\
\midrule
T1 — valeur-seuil enfouie & Constante critique sans nom ni justification & constante
nommée + rationale dans la spec I \\
T2 — convention d’interface implicite & Structure ou ABI déclarée de mémoire & contrat
typé + invariant de taille testé \\
T3 — dépendance à l’hôte & Comportement dépendant de l’OS d’exécution & contrat
indépendant de l’hôte, union fixée \\
T4 — parsing de surface fragile & Analyse de libellés localisés ou volatils & fragments
indépendants des libellés + fixtures partagées \\
T5 — contrat dupliqué & Même connaissance codée en N exemplaires divergents & source
unique : spec I + fixtures \\
T6 — garantie non testée & Comportement critique sans vecteur de rejet & cas d’échec
nommé + test de rejet \\
\bottomrule
\end{tabularx}
\end{table}

Deux règles d’usage : tout défaut corrigé est classé dans cette taxonomie au journal des
décisions (la métrique « défauts ouverts par classe » du tableau de bord en dérive) ; la
taxonomie sert de liste de contrôle de rédaction — chaque bloc du gabarit verrouille au
moins une classe. Les occurrences de référence viennent de la démonstration : le seuil de
détection non justifié (T1+T6), la structure \code{SP\_DEVINFO\_DATA} de 424 octets au
lieu de 32 (T2, échec \code{0x6F8}), la sanitisation de noms dépendante de l’hôte (T3), le
parsing de pnputil sur libellés français (T4+T5, trois copies).

\section{Les trois couches : architecture et traçabilité}\label{sec:couches}

Le diff d’intention descend de I vers C, le verrou croisé relie C et B dans les deux sens,
et la couche B seule quitte le laboratoire. \textbf{Aucune couche n’est optionnelle} —
c’est une propriété normative : privé de I, le système perd son ancrage sémantique ; privé
de C, l’agent devient aveugle ; privé de B, il n’y a plus rien à exécuter. Toute
implémentation qui livrerait B sans I et C à jour viole R1 ; toute évolution qui
contournerait C pour éditer B directement viole R3 et la boucle.

\begin{table}[htbp]
\centering\small
\caption{Les trois couches.}
\label{tab:couches}
\begin{tabularx}{\textwidth}{@{}l l Y Y@{}}
\toprule
\textbf{Couche} & \textbf{Visibilité} & \textbf{Contenu} & \textbf{Rôle} \\
\midrule
I & Privée & Spécifications normalisées : contrats pré/post, invariants, cas d’échec,
vecteurs & Compression de l’intention \\
C & Privée, signée, chiffrée & Index des blocs, adresses, preuves, manifestes,
traçabilité & Unique interface de développement \\
B & Livrée & Code machine scellé, padding, checksums runtime & Exécution pure,
sémantiquement muette \\
\bottomrule
\end{tabularx}
\end{table}

La séparation croise une seconde dimension : la cible. La couche I est indépendante de la
cible — un contrat s’écrit une fois, ses types abstraits étant paramétrés par le profil.
La couche B est par profil. La couche C indexe l’ensemble : chaque entrée porte son
profil, et frag-bld organise les empreintes par profil. Quatre invariants structurels,
vérifiés à chaque cycle — leur violation est un refus de build : B n’est jamais nu (R1) ;
les adresses exposées sont stables (R7, par profil) ; la signature couvre la carte entière
(R9) ; B est auto-vérifiant au chargement (R10 renforcée).

\subsection{Traçabilité I $\leftrightarrow$ C $\leftrightarrow$ B $\leftrightarrow$
tests}\label{sec:traceabilite}

Chaque entrée de la carte référence quatre maillons : la spécification (identifiant +
révision), l’entrée d’index et le contrat embarqué, le(s) bloc(s) instancié(s) avec
empreinte(s) par profil, et les vecteurs exécutés avec le verdict du dernier cycle. Trois
propriétés sont contrôlées mécaniquement par la \textbf{commande d’audit de traçabilité} :
aucune spécification sans entrée de carte (l’exception visible étant la reverse-spec en
statut « prête ») ; aucun bloc scellé sans spécification référencée ; aucun vecteur
orphelin. Une chaîne incomplète est un refus de build — au même titre qu’un binaire nu.
La traçabilité est la condition de l’audit par des tiers : elle permet de répondre à
« quelle exigence ce bloc implémente-t-il, et quel test la défend ? » sans lecture du
binaire.

\section{La carte : fragmentation et budgets}\label{sec:carte}

\begin{sloppypar}
La carte est le cœur opérationnel du protocole : elle rend le binaire navigable, porte les
preuves et concentre la valeur intellectuelle. Elle est autodescriptive : un en-tête
normalisé consigne la version du protocole, la suite cryptographique~— hachage, signature,
chiffrement, chacun dans une liste blanche (D15)~— et les profils actifs. Sa spécification
détaillée fait l’objet de \code{FORMAT\_CARTE.md} — fragments, format d’entrée, verrou,
cycle de vie, export, re-cryptage —, mise à jour avec la présente révision : l’en-tête
porte \code{version=0.4.0} et chaque entrée de bloc porte le champ de traçabilité
\code{spec=<DOM>-<NN>:R}.
\end{sloppypar}

\paragraph{Norme de fragmentation.} La carte est fragmentée en quatre — frag-loc, frag-con,
frag-prf, frag-bld — et la v0.4.0 pose la norme : aucune opération de la boucle n’exige la
lecture de plus d’un fragment ; \textbf{aucune opération n’exige la lecture intégrale du
binaire} — le balayage est interdit (§\ref{sec:boucle}) ; l’audit complet de la carte est
une opération exceptionnelle, non budgétée, journalisée.

\paragraph{Budgets par type d’opération.}

\begin{table}[htbp]
\centering\small
\caption{Budgets en tokens par opération (localiser amendé à 64 tokens par D31 après la mesure M1 — le gabarit canonique v0.4.0 mesurait 43 tokens en cl100k\_base ; autres budgets inchangés depuis la v0.3).}
\label{tab:budgets}
\begin{tabularx}{\textwidth}{@{}Y Y l@{}}
\toprule
\textbf{Opération} & \textbf{Données lues} & \textbf{Budget} \\
\midrule
Localiser un bloc & frag-loc — une entrée & 64 tokens max (D31) \\
Éditer simplement un bloc & frag-loc + frag-con + notation & 3 000 tokens max \\
Ajouter un bloc & frag-loc + frag-con + seuils frag-prf & 6 000 tokens max \\
Instancier un profil nouveau & I partagée + frag-con + notation du profil & 6 000 tokens
max par bloc \\
Audit complet (rare) & carte entière déchiffrée & non budgété \\
\bottomrule
\end{tabularx}
\end{table}

Indicateurs de santé : précision de localisation (cible 95 \% dès M0), coût moyen d’une
édition simple (cible 5 000 tokens), taux de régression post-patch (cible zéro).
Instrumentés depuis M1 : la mesure effective (51 / 719 / 1 088 tokens, D31) a corrigé
le budget de localisation à 64 et valide les deux autres ; le tableau de bord les
affiche comme métriques mesurées, avec leur écart.
